\section{Weighted forests and complete flags}
Let \(T\) be a finite tree with vertex set \(V\) and edge set \(E\). Give each vertex a finite nonempty block \(B_v\), with the blocks pairwise disjoint, and put \(\Omega=\bigsqcup_vB_v\). For \(F\subseteq E\), let \begin{equation}
K_F=\prod_C S_{\bigcup_{v\in C}B_v},
\label{eq:03-forest:2-1}
\end{equation} where \(C\) runs over the components of the forest \((V,F)\). Root \(T\), and let \(A_e\subseteq\Omega\) be the union of blocks on the side of \(T-e\) away from the root. In the finite symmetric group \(S_\Omega\), set \(H_e=K_{E\setminus\{e\}}=S_{A_e}\times S_{\Omega\setminus A_e}\).

\textbf{Theorem UB.F (forest intersections and flag completion).} The following hold.

\begin{enumerate}
\def\labelenumi{\arabic{enumi}.}
\tightlist
\item
  \(K_F\cap K_D=K_{F\cap D}\). Moreover \(K_F\) is generated by \(K_\varnothing\) together with the groups \(K_{\{e\}}\), \(e\in F\).
\item
  The coset system of the \(H_e\) is a firm, residually connected, flag-transitive geometry. Its type-\(e\) elements identify with the \(|A_e|\)-subsets of \(\Omega\), and distinct types satisfy \begin{equation}
  X_e\mathrel{*}X_f\quad\Longleftrightarrow\quad
  |X_e\cap X_f|=|A_e\cap A_f|.
  \label{eq:03-forest:2-2}
  \end{equation}
\item
  Its Borel subgroup is \(K_\varnothing=\prod_vS_{B_v}\). The residue whose remaining types are \(F\) is the join of the same constructions on the nontrivial components of \((V,F)\). Its acting group is \(K_F\); components with no edges contribute only a kernel on this residue.
\end{enumerate}

\emph{Proof.} A permutation belongs to both groups in (1) exactly when it preserves every intersection of their component blocks. Two vertices lie in a common block of both forests precisely when their unique tree path belongs to both edge sets. Thus the common refinement is the component partition for \(F\cap D\). Transpositions inside the blocks, together with transpositions joining adjacent blocks along retained edges, connect each component; they generate its full symmetric group. This proves (1).

The coset \(gH_e\) maps to \(gA_e\). Two pairs of subsets are conjugate under \(S_\Omega\) exactly when their four Boolean atoms have the same sizes, giving \eqref{eq:03-forest:2-2}. The sets \(A_e\) are laminar: any two are nested or disjoint. Equation \eqref{eq:03-forest:2-2} forces the identical nesting and disjointness relations on any pairwise incident family \((X_e)\). Add \(\Omega\) as a top container. From each selected container subtract its maximal selected proper subcontainers. These differences partition \(\Omega\), and the corresponding differences for the \(A_e\) and \(X_e\) have equal sizes. Bijections between these atoms combine into a single permutation sending every selected \(A_e\) to \(X_e\). This proves (G3) for every partial flag.

Repeated use of (1) gives \(H_J=K_{E\setminus J}\). Different edge sets give different component partitions and different groups, since every block is nonempty and adding an edge allows a transposition between its two sides. Thus (G1) holds. If at least two edge types remain, the groups obtained by deleting each remaining edge collectively contain \(K_\varnothing\) and each remaining one-edge group, proving (G2). The coset criterion gives the geometry assertions. Its residue group and its subgroups factor over the forest components; all elements from different components are incident, which gives the stated join and the kernel convention. ∎

For unit blocks \(B_v=\{v\}\), write \(\tau_e\) for the endpoint transposition of \(e\). The theorem says \begin{equation}
K_F=\langle\tau_e:e\in F\rangle
 =\prod_C S_C,\qquad K_\varnothing=1.
\label{eq:03-forest:2-3}
\end{equation} This is the unweighted model used for {[}BC98{]}. The weighted formulation also supplies the parity construction after an edge contraction and the disjoint-subset construction after a star specialization.
